Summary: The authors present a video classifier for Indian classical dance that relies on a Multi-Frame Multi-Head self-Attention Layer (MFMHAL) for any CNN backbones. Specifically, they implement MFMHAL on ResNet50. To support this study, the authors construct an Indian classical dance (ICD) video dataset, which includes 5 songs with 10 performances. The experiments on the proposed ICD dataset evaluate the performance of their approach and analyze the training parameters for ablation studies. The results show that classification accuracy improves by over 8% across all datasets.

Reproducibility: The implementation details of this paper support the reproducibility of this work, and the open-source of their own dataset can further aid in this regard.

Experiments:

The comparisons of the architectures between MFMHLA and other baselines, including [37], could benefit from a comparison with the Transformer-only module for backbone evaluation.

While the backbone for most experiments is ResNet-50, an ablation study on different backbone nets, such as ResNet-34 or 101, could further evaluate the improvements of MFMHAL.

The comparisons in Fig. 8 with ResNet-34 seem unfair, as the proposed work is based on ResNet-50.

Tables/Figures:

For most figures, such as Fig. 8/9, the notations are too small to follow. These figures need improvement by increasing the font size.

In Tab.3, the values of "Their Own Dance Dataset" are not clear.

More visualizations from other datasets, such as Fig. 7, can support the generalization of this MFMHAL design.










Here's my suggested revision:

Support from experimental results:

Fig. 7 illustrates that MFMHAL focuses on the skin, such as hands and feet, rather than clothing, while the results in Fig. 8 convincingly demonstrate the improvement in mAR.

Tab. 8 and Tab. 9 show the quantitative results with many well-known state-of-the-art attention mechanisms, demonstrating the effectiveness of the proposed approach.

Contributions:

The contributions of this paper are clear and convincing. Firstly, the authors propose a Multi-Frame Multi-Head self-Attention Layer (MFMHAL) for convolutional neural layers. Secondly, they introduce an Indian Classical Dance dataset to support their specific task. Finally, they evaluate MFMHAL on both the proposed ICD dataset and other benchmarks by comparing it with other state-of-the-art methods.

Limitations:

The authors do not specifically address their limitations. It would be helpful to elaborate on this in the experiments or conclusion section.

It appears that the MFMHAL module involves more parameters than the original ResNet backbone, which may increase the computational resources required and inference time.

Writing:

On page 2, line 17, please correct the spacing of "multi-frame multi-head self-attention."

Please add a limitations section to further discuss the potential disadvantages of this work.












我希望你能担任英语拼写校对和修辞改进的角色, 用于专业的论文评审工作。我会用英文语言和你交流，你会识别语言，将其润色为优美和精炼的英语。请将我简单的词汇和句子替换成更为优美和高雅的表达方式，确保意思不变，但使其更具文学性。请回答更正和改进的部分，并进行解释。

Summary: The authors present a video classifier of Indian classical dance, which relied on a Multi-Frame Multi-Head self-Attention Layer (MFMHAL) for any CNN backbones. Specifically, they implement MFMHAL on the ResNet-50. To support this study, the authors construct an Indian classical dance (ICD) video dataset, which includes 5 songs with 10 performances. The experiments on the proposed ICD dataset evaluate the performance of their approach and analyze the training parameters for ablation studies. The results show that classification accuracy improves by over 8% across all datasets.

Reproducibility: The implementation details of this paper support the reproducibility of this work, while the open-source of their own dataset can further help.

Experiments: 
1. The comparisons of the architectures between the MFMHLA and other baselines are extensive. Including [37], it could help to compare with the Transformer-only module for the backbone evaluation.

2. For most experiments, the backbone is ResNet-50. While the ablation study on different backbone nets, like ResNet-34 or 101, can further evaluate the improvements of the MFMHAL.

3. The comparisons in Fig. 8 with the ResNet-34 seem to be unfair, while the proposed work is based on ResNet-50.

Tables/Figures: 
1. For most figures, like Fig. 8/9, the notations are too small to follow. These figures still need to improve by increasing the font size.

2. Tab.3 The values of "Their Own Dance Dataset" are not clear to me.

3. More visualizations from other datasets, like Fig. 7, can support the generalization of this MFMHAL design.

The support from experimental results:
1. Fig. 7 illustrates that MFMHLA focuses on the skins, like hands and feet, rather than clothing, while the results in Fig. 8 are convincing to show their improvement on mAR.

2. Tab. 8 and Tab. 9 show the quantitative results with many well-known state-of-the-art attention mechanisms. The results show the effectiveness of the proposed works.

Contributions:
The contributions of this paper are clear and convincing. 1) The authors propose a Multi-Frame Multi-Head self-Attention Layer (MFMHAL) for convolutional neural layers. 2) They propose an Indian Classical Dance dataset to support their specific task. 3) They evaluate MFMHAL on both the proposed ICD dataset and other benchmarks by comparing it with other state-of-the-art.

Limitations: 
1. The authors do not introduce their limitations specifically. Please elaborate on this in your experiments or conclusion section.

2. It seems that the MFMHAL module involves more parameters than the original ResNet backbone. Therefore, it costs more computer resources and increases the inference time. 

Writings: 
1. Page 2 - Line 17 - Right: Please correct the space of "multi-framemulti-headself-attention".
2. Please add the limitation section to further introduce the disadvantages of this work.






DINAR, an approach for generating a neural human avatar from a single front image. DINAR consists of two components: an overall avatar neural rendering pipeline and a diffusion-based inpainting model on a 2D human-aware texture domain. The former is performed in an end-to-end fashion, which includes a StyleGAN2 module for generating neural texture and a neural renderer for the final avatar rendering. The latter is performed on the latent space of VQGAN to inpaint the missing information from the front view. The results presented in the paper and video are , particularly for this challenging task.

strength:
1. The proposed approach is technically sound, incorporating diffusion on the UV latent space and neural avatar rendering to generate satisfactory neural avatar results.

2. According to the authors, this work could be the first to extend diffusion models to the task of generating human body textures and outperform many existing state-of-the-art (SOTA) results in a one-shot neural avatar setting.

weakness:
1. Writing and Illustration: a) In line 463, do the authors try to explain how to generate GT texture in the Texture Merging Part (Sec3.2)? Or how their work extends from one-shot to few-shot? According to line 538, it is for GT generation. The arrangement is quite confusing. Please elaborate on this before this subsection or make it into a new training subsection. b) Fig.2 seems to only explain the avatar part. The relation between this and the diffusion part is not that obvious. Does the approach try to replace the partial texture from the RGB sampler with the in-painted texture? Or should the diffusion part be processed in the middle of this VQGAN? An illustration of the inpainting results on Fig.2 could help. c) The paper and supplemental materials provide many details and figures, but there is no inference pipeline, which causes confusion in distinguishing between training and inference.

2. The motivation behind StyleGAN2: The motivation of this part in the whole pipeline is not convincing to me. I understand that the goal is to provide more global and local information of appearance from the input front image directly. However, why use this module? Can this module significantly improve the quality? An evaluation and discussion should be provided.

3. Insufficient ablation: Ablation should play an important role in your work, especially since there are many modules in the whole pipeline. However, there is only one figure without dynamic results of ablation evaluation, without any quantitative results, and even without the back view in Fig.7. This is quite insufficient and weak. If the reviewer ignored any ablation study, please indicate.

Rating:
I think that this work is technically sound, and the quality of avatar rendering results is satisfactory with the challenging input of only one single image. However, the authors should improve on the ablation study, reorganize the approach section, and elaborate on the motivation and improvement for the StyleGAN2 part. Due to the above concerns, which are not clear to me at the moment, I would give a borderline rating now and would modify the rating according to the rebuttal and discussion afterward.





 
 % 1 background - large model and multi-modal
  With the development of pretrained large language models, a uniform model for language and other multimodal, like motion, is emerging.
  %
  % 2 Insight/observation - motion task small task,
  %   a uniform pre-trained motion framework, which can extend many motion-relevant task.
  % The realm of human motion research suffers from xxx challenges and hurdles. 
  % 
  % 3 Propose/motivation - motion as body language, semantic information as words in the language
  Fortunately, human motion exhibits a coupled structure similar to human language, which can be construed as a body language form. 
  % 
  % 4 method - propose MotionGPT
  % text and motion in a unified manner.
  Motivated by this, we propose MotionGPT, a unified motion-language model that achieves state-of-the-art performance on motion-language tasks.
  %     
  % 5 details - specifically, motion quaternization xx, language backbone xx
  Specifically, we advance the discrete vector quantization for human motion and transfer 3D motion as motion tokens, similar to word embeddings. 
  %
  Based on this "motion vocabulary", we perform language modeling on both motion and text in a unified manner, treating human motion as a specific language.
  % 
  Inspired by the prompts learning, we pre-train MotionGPT with the mix of motion-language data and finetune it with prompt-based motion QA tasks.
  %
  % 6 Experimental results - xx
  Extensive experiments demonstrate that MotionGPT achieves state-of-the-art performance on motion generation, motion caption, motion prediction, and motion QA tasks.

  
% MLD
% We study a challenging task, conditional human motion generation, which produces plausible human motion sequences according to various conditional inputs, such as action classes or textual descriptors. Since human motions are highly diverse and have a property of quite different distribution from conditional modalities, such as textual descriptors in natural languages, it is hard to learn a probabilistic mapping from the desired conditional modality to the human motion sequences. Besides, the raw motion data from the motion capture system might be redundant in sequences and contain noises; directly modeling the joint distribution over the raw motion sequences and conditional modalities would need a heavy computational overhead and might result in artifacts introduced by the captured noises. To learn a better representation of the various human motion sequences, we first design a powerful Variational AutoEncoder (VAE) and arrive at a representative and low-dimensional latent code for a human motion sequence. Then, instead of using a diffusion model to establish the connections between the raw motion sequences and the conditional inputs, we perform a diffusion process on the motion latent space. Our proposed Motion Latent-based Diffusion model (MLD) could produce vivid motion sequences conforming to the given conditional inputs and substantially reduce the computational overhead in both the training and inference stages. Extensive experiments on various human motion generation tasks demonstrate that our MLD achieves significant improvements over the state-of-the-art methods among extensive human motion generation tasks, with two orders of magnitude faster than previous diffusion models on raw motion sequences.


 From Image as a Foreign Language: BEIT Pretraining for All Vision and Vision-Language Tasks
 1 A big convergence of language, vision, and multimodal pretraining is emerging. 
 2 In this work, we introduce a general-purpose multimodal foundation model BEIT-3, which achieves state-of-the-art transfer performance on both vision and visionlanguage tasks. 
 3 Specifically, we advance the big convergence from three aspects: backbone architecture, pretraining task, and model scaling up. 
 4 We introduce Multiway Transformers for general-purpose modeling, where the modular architecture enables both deep fusion and modality-specific encoding. 
 5 Based on the shared backbone, we perform masked “language” modeling on images (Imglish), texts (English), and image-text pairs (“parallel sentences”) in a unified manner. 
 6 Experimental results show that BEIT-3 obtains state-of-the-art performance on object detection (COCO), semantic segmentation (ADE20K), image classification (ImageNet), visual reasoning (NLVR2), visual question answering (VQAv2), image captioning (COCO), and cross-modal retrieval (Flickr30K, COCO).


% Challenging Task: Conditional 3D shapes generation. 
% 1. Domain gap (theoretically); 2. lack of data (practical)
Generating realistic 3D shapes conditioned on 2D images or texts is challenging due to vast domain gap between shape-image-text modalities and different 3D objects have highly variant topology structure.%, where the absence of large-scale 3D dataset further intensify the difficulty of this task.
% Contrast Model aligns multiple modalities.
Fortunately, recent advances of contrastive model show strong ability for extracting visual and textual information into an aligned embedding space, which also boost the text-to-image generation.
% Why we wanna align three modalities.
Motivated by this, representing 3D shapes into an aligned space could be promising to close the domain gap among three modalities and further facilitate multi-modal-conditioned 3D shape generation.
% Our model
Therefore, we devise a two-stage model to learn an aligned representation among multiple modalities information and generate high-quality 3D shapes based on learnt align space.
% Aligned Shape VAE
The first-stage-model is a shape-image-text aligned variational auto-encoder (SITA). It represents 3D shapes as shape tokens aligned to image embedding and text embedding. and reconstructs fine-grained neural fields corresponding to given shape tokens.
% CDM
With aligned representation, the second-stage-model, conditional align-latent diffusion model (CAL) generates shape tokens given images or texts.
% New Metric
Moreover, the aligned embedding space enables directly compute similarities between 3D shapes and conditional information without any projection, which provides a valid metric for evaluation. 
% EXP
We evaluate the proposed framework on shape reconstruction, image-conditional and text-conditional 3D shape generation, shape classification and shape retrieve tasks. Extensive experiments validate that our model could generate high-quality 3D shapes confirm to conditional information. 











# introduction 

% language
Recent years have witnessed the rapid development of increasingly sophisticated pre-trained language models like GPT and BERT, which have demonstrated remarkable performance in various natural language processing tasks. 
% motion-language
However, there has been a lack of research on utilizing language models for human motion-related tasks, despite the potential benefits in fields such as gaming, robotics, and human behavior analysis. 
This gap in research highlights the need for a framework that can effectively leverage the strong language generation and zero-shot transfer abilities of pre-trained language models for human motion-related tasks.
% Human motion holds great potential in a wide range of applications.
% For instance, animating virtual characters that move naturally and realistically is essential in the gaming and entertainment industry. Moreover, robots can benefit from understanding human motion as a means of improving their interaction with humans in various settings. In addition, human motion analysis can provide insights into users' behavior, which is particularly useful in areas such as healthcare and sports
